\documentclass[12pt]{article}
\usepackage[utf8]{inputenc}
\usepackage[T1]{fontenc}
\usepackage[margin=1in]{geometry}
\usepackage{setspace}
\usepackage{xcolor}
\usepackage[colorlinks=true, urlcolor=blue!60!black, linkcolor=blue!60!black]{hyperref}
\setlength{\parindent}{2em}
\setlength{\parskip}{1.25em}
\renewcommand{\baselinestretch}{1.0}
\begin{document}
\begin{center}
{\LARGE \textit{The Civil Rights Cases}, 109 U.S. 3 (1883)\par}
\vspace{0.8em}
{\large Case Brief --- Four-Page Edition\par}
\end{center}
\vspace{1em}

\section{Memory Jogger}

In an 8--1 decision establishing the modern state action doctrine, the Court struck down the public accommodations provisions of the Civil Rights Act of 1875, holding that the Fourteenth Amendment restricts only state action and that the Thirteenth Amendment does not reach private racial discrimination in inns, theaters, and railroads (\href{https://www.courtlistener.com/opinion/90897/civil-rights-cases/}{CourtListener}).

\section{Detailed Case Facts}

Section 1 of the Civil Rights Act of 1875 entitled all persons to ``the full and equal enjoyment of the accommodations, advantages, facilities, and privileges of inns, public conveyances on land or water, theatres, and other places of public amusement,'' regardless of race or previous condition of servitude. Section 2 made any violation by any person a federal misdemeanor and created a \$500 civil penalty for the aggrieved party.

Five disputes were consolidated: indictments in Kansas and Missouri against innkeepers who denied Black guests accommodations; a California information against a man who refused a Black patron a seat in the dress circle of Maguire's Theatre in San Francisco; a New York indictment over the Grand Opera House; and a civil suit by Richard and Sallie Robinson after a conductor excluded Mrs. Robinson, a Black woman holding a ticket from Tennessee to Virginia, from the ladies' car.

\section{Procedural History}

In the Kansas, Missouri, and New York cases, the two circuit judges divided on the Act's constitutionality and certified the question. In California the circuit court sustained a demurrer, and the United States sought review by writ of error. In Tennessee the jury returned a verdict for the railroad, and the Robinsons sought review. The Supreme Court decided all five cases on October 15, 1883.

\section{Judicial Votes}

Sections 1 and 2 held void, \textbf{8--1}.
\begin{itemize}
\item \textbf{Majority:} Justice Bradley, joined by Chief Justice Waite and Justices Miller, Field, Woods, Matthews, Gray, and Blatchford.
\item \textbf{Concurrences:} None.
\item \textbf{Dissent:} Justice Harlan.
\end{itemize}

\section{Holding}

The Fourteenth Amendment prohibits only state action; Section 5 empowers Congress to enact \textit{corrective} legislation against state laws and official acts, not direct legislation regulating private conduct. The Thirteenth Amendment authorizes direct legislation against slavery's ``badges and incidents,'' but a private proprietor's denial of equal accommodations is not one. Bradley reserved whether access is a right against state interference, whether the Act was valid in the Territories and the District of Columbia, and whether Congress could regulate interstate conveyances under the Commerce Clause.

\section{Analysis of Opinions}

\textbf{Majority (Bradley).} ``No State shall'' targets ``State action of a particular character.'' Like the Contracts Clause, a prohibition on state wrongdoing supports only remedies predicated on state laws, customs with the force of law, or official conduct. Because the 1875 Act reached private individuals directly, even in States with fair laws, it created a federal ``code of municipal law'' in violation of the Tenth Amendment. Bradley contrasted Section 4 of the Act, upheld in \textit{Ex parte Virginia} because state officers excluded jurors, and the 1866 Act's ``under color of'' state law provision. The Thirteenth Amendment reaches the inseparable incidents of slavery, such as inability to contract, sue, or hold property, but treating every private refusal of service as a badge of slavery would be ``running the slavery argument into the ground.''

\textbf{Harlan (dissenting).} The majority's grounds were ``entirely too narrow and artificial.'' Under \textit{Prigg} and the Fugitive Slave Acts, the Court had implied direct federal power to protect slaveholders, yet now denied express power to protect freedom. Because slavery rested on race, freedom entails immunity from racial discrimination in civil rights. Railroads, common-law inns, and licensed theaters perform public functions and act as ``agents or instrumentalities of the State.'' The Citizenship Clause is an affirmative grant Congress may protect directly.

\section{Examples: Future Applications}

\begin{enumerate}
\item \textbf{Same-Side (Struck Down): Private Violence Remedy.} A federal civil remedy against private perpetrators of gender-motivated violence exceeds Section 5, as in \textit{United States v. Morrison} (2000).
\item \textbf{Same-Side (No Claim): Private Club.} A members-only club with no state funding or franchise that excludes an applicant by race commits ``individual invasion of individual rights,'' not state action.
\item \textbf{Opposite-Side (State Action): Covenant Enforcement.} A state-court injunction enforcing a racial covenant supplies state action, as \textit{Shelley v. Kraemer} (1948) held.
\item \textbf{Opposite-Side (Valid): Other Powers.} Public accommodations laws under the Commerce Clause (\textit{Heart of Atlanta Motel}, 1964) and housing laws under the Thirteenth Amendment (\textit{Jones v. Alfred H. Mayer}, 1968) are valid.
\item \textbf{Fence-Sitter: Leased Public Space.} A segregating restaurant in a city-owned garage is private under Bradley's reading but entwined with the State under Harlan's view.
\end{enumerate}

\section{Critique}

\textbf{Originalist views.} Defenders stress the text ``No State shall'' and the Tenth Amendment, and Richard Epstein notes the doctrine protects a private sphere. Michael McConnell shows that Reconstruction Congresses read the Amendment more robustly; Randy Barnett and Evan Bernick argue that equal protection obligated States to protect against private wrongs; and Pamela Brandwein contends Bradley preserved a ``state neglect'' concept later erased.

\textbf{Progressive views.} Scholars place the case in the anti-canon with \textit{Plessy}. Charles Black argued that the public/private line is illusory because exclusion rests on state trespass law, and critics fault the Court for forcing Congress to use the Commerce Clause, an economic power, to address a moral wrong.

\section{Key Quotations}

\begin{quote}
``It is State action of a particular character that is prohibited. Individual invasion of individual rights is not the subject matter of the amendment.''
\end{quote}
\noindent --- Justice Bradley, Opinion of the Court, 109 U.S. at 11.

\begin{quote}
``In fine, the legislation which Congress is authorized to adopt in this behalf is not general legislation upon the rights of the citizen, but corrective legislation, that is, such as may be necessary and proper for counteracting such laws as the States may adopt or enforce, and which, by the amendment, they are prohibited from making or enforcing, or such acts and proceedings as the States may commit or take, and which, by the amendment, they are prohibited from committing or taking.''
\end{quote}
\noindent --- Justice Bradley, Opinion of the Court, 109 U.S. at 13--14.

\begin{quote}
``When a man has emerged from slavery, and, by the aid of beneficent legislation, has shaken off the inseparable concomitants of that state, there must be some stage in the progress of his elevation when he takes the rank of a mere citizen and ceases to be the special favorite of the laws, and when his rights as a citizen or a man are to be protected in the ordinary modes by which other men's rights are protected.''
\end{quote}
\noindent --- Justice Bradley, Opinion of the Court, 109 U.S. at 25.

\end{document}
